\chapter{МАТЕМАТИЧЕСКАЯ МОДЕЛЬ И ЧИСЛЕННЫЕ МЕТОДЫ}

В главе формализуется математическая постановка расчёта
термодинамических свойств воды и водяного пара.
Нормативной основой служит IAPWS-IF97 \cite{iapws_if97_2012}.
Стандарт задаёт региональную топологию, уравнения состояния и часть обратных зависимостей.
Поэтому вычислительный алгоритм включает обязательные шаги
классификации области, выбора модели и контроля применимости.
Принятые здесь соотношения и численные процедуры
непосредственно используются в реализации \texttt{if97\_core},
описанной в следующей главе.

\section{Разбиение области IAPWS-IF97 на регионы}

Рабочая область IF97 разделена на регионы,
в каждом из которых применяются собственные уравнения
и собственные нормировки безразмерных переменных.
Такое разбиение отражает физические различия
между сжатой жидкостью, перегретым паром и околокритической областью.
Одна универсальная аналитическая формула для всей области
в промышленной постановке не обеспечивает одновременно
точность, устойчивость и вычислительную эффективность.

Следовательно, для любого маршрута расчёта
определение региона является обязательным этапом.
Только после классификации состояния корректно
выбирается модель и вычисляются итоговые свойства.

\subsection{Регионы 1--5}

IF97 выделяет пять основных регионов и отдельную область насыщения.
Регионы различаются физическим смыслом
и используемым термодинамическим потенциалом.
Сводная характеристика приведена в таблице~\ref{tab:if97_regions_overview}.

\begin{table}[h]
\centering
\caption{Регионы IAPWS-IF97 и используемые формулировки}
\label{tab:if97_regions_overview}
\begin{tblr}{
  width=\linewidth,
  colspec={|X[1,l]|X[2,l]|X[2,l]|},
  row{1}={font=\bfseries},
}
\hline
Регион & Физический смысл & Потенциал и независимые переменные \\
\hline
1 & Сжатая вода (подохлаженная жидкость) & Гиббс, $p$ и $T$ \\
\hline
2 & Перегретый пар & Гиббс, $p$ и $T$ \\
\hline
3 & Высокоплотная околокритическая область & Гельмгольц, $\rho$ и $T$ \\
\hline
4 & Линия насыщения и двухфазная область & Насыщение, $p$ или $T$ и степень сухости $x$ \\
\hline
5 & Высокотемпературная область (пар) & Гиббс, $p$ и $T$ \\
\hline
\end{tblr}
\end{table}

В инженерных сценариях вход задаётся парами параметров,
доступных из измерений или расчётных схем.
В работе используются параметры $p$, $T$, $h$, $s$, $\rho$, $x$.
На выходе формируется единый набор свойств:
$v$, $\rho$, $h$, $s$, $u$, $c_p$, $w$ и номер региона IF97.

\subsection{Линия насыщения и граница B23}

Линия насыщения (Region~4) является фазовой границей
между жидкостью и паром.
На этой линии давление и температура связаны однозначно.
Состояние задаётся насыщенными свойствами ($x=0$ или $x=1$)
или двухфазной смесью при $0<x<1$.
Внутри двухфазной области давление определяется температурой насыщения,
а параметр $x$ задаёт массовую долю пара.

Переходы через Region~4 требуют отдельного контроля,
поскольку производные свойств изменяются резко.
Для двухфазной области характерны
провал скорости звука и рост изобарной теплоёмкости.
Эти эффекты учитываются при выборе численных процедур
и при интерпретации диаграмм.

Граница B23 отделяет Region~2 от Region~3.
Она определяет момент перехода
от формулировки через потенциал Гиббса
к формулировке через потенциал Гельмгольца.
Корректная обработка B23 обязательна
для маршрутизации в постановках $(p,T)$, $(p,h)$ и $(p,s)$.

\section{Термодинамические потенциалы и расчётные зависимости}

Модель IF97 вычисляет свойства через производные
безразмерных термодинамических потенциалов.
Такое представление обеспечивает единую форму записи,
корректное масштабирование переменных
и лучшую численную устойчивость.

В регионах 1, 2 и 5 используется потенциал Гиббса.
В регионе 3 используется потенциал Гельмгольца.
Region~4 описывается зависимостями насыщения и правилами смешения.

\subsection{Формулировка через потенциал Гиббса для регионов 1, 2 и 5}

Для регионов 1, 2 и 5 стандарт IF97 задаёт безразмерную энергию Гиббса
\begin{equation}
\gamma(\pi, \tau) = \frac{g(p, T)}{R T},
\end{equation}
где $g$ --- удельная энергия Гиббса, $R$ --- газовая постоянная воды и водяного пара,
а безразмерные переменные определяются как
\begin{equation}
\pi = \frac{p}{p^{*}}, \qquad \tau = \frac{T^{*}}{T}.
\end{equation}
Нормировочные константы $p^{*}$ и $T^{*}$ задаются стандартом отдельно для каждого региона.

Свойства выражаются через частные производные $\gamma(\pi,\tau)$.
Используются обозначения
$\gamma_{\pi} = \partial \gamma/\partial \pi$,
$\gamma_{\tau} = \partial \gamma/\partial \tau$,
$\gamma_{\pi\pi} = \partial^{2}\gamma/\partial \pi^{2}$,
$\gamma_{\tau\tau} = \partial^{2}\gamma/\partial \tau^{2}$.

Ключевые зависимости, используемые в вычислениях, включают:
\begin{equation}
v = \frac{R T}{p}\,\pi\,\gamma_{\pi},
\end{equation}
\begin{equation}
h = R T\,\tau\,\gamma_{\tau},
\end{equation}
\begin{equation}
s = R \left(\tau\,\gamma_{\tau} - \gamma\right),
\end{equation}
\begin{equation}
u = R T \left(\tau\,\gamma_{\tau} - \pi\,\gamma_{\pi}\right),
\end{equation}
\begin{equation}
c_{p} = -R\,\tau^{2}\,\gamma_{\tau\tau}.
\end{equation}

Скорость звука $w$ и остальные производные свойства
вычисляются по стандартным комбинациям производных $\gamma$
\cite{iapws_if97_2012}.

\subsection{Формулировка через потенциал Гельмгольца для региона 3}

Для Region~3 стандарт использует формулировку через безразмерную энергию Гельмгольца
\begin{equation}
\phi(\delta, \tau) = \frac{f(\rho, T)}{R T},
\end{equation}
где $f$ --- удельная энергия Гельмгольца,
а безразмерные переменные определяются как
\begin{equation}
\delta = \frac{\rho}{\rho^{*}}, \qquad \tau = \frac{T^{*}}{T}.
\end{equation}
Нормировочные константы $\rho^{*}$ и $T^{*}$ задаются стандартом.

В этой формулировке давление выражается через производную $\phi$ по плотности:
\begin{equation}
p = \rho R T\,\delta\,\phi_{\delta}.
\end{equation}
Это соотношение показывает,
что для Region~3 естественными независимыми переменными
являются $\rho$ и $T$.
При постановке $(p,T)$ требуется сначала определить плотность,
согласованную с заданным давлением.

После определения $\rho$ остальные свойства могут быть вычислены через производные $\phi$.
Например, энтропия и энтальпия выражаются как:
\begin{equation}
s = R \left(\tau\,\phi_{\tau} - \phi\right),
\end{equation}
\begin{equation}
h = R T \left(\tau\,\phi_{\tau} + \delta\,\phi_{\delta}\right).
\end{equation}
Остальные свойства (включая $u$, $c_p$ и $w$)
получаются по формулам IF97 через производные $\phi$
до второго порядка \cite{iapws_if97_2012}.

\section{Прямые и обратные маршруты расчёта}

Практическая применимость модели определяется набором маршрутов расчёта.
Маршрут задаётся входной парой параметров
и определяет последовательность вычислительных шагов.
В проекте используются как прямые зависимости IF97,
так и обратные постановки с аналитическими и численными процедурами.

\subsection{Прямой маршрут pT}

Прямой маршрут $pT$ является базовым для IF97.
По заданным давлению $p$ и температуре $T$ выполняются следующие этапы:
\begin{enumerate}
\item валидация диапазонов входных значений в пределах области применимости стандарта;
\item определение региона IF97 по топологическим границам (линия насыщения, граница B23 и др.);
\item вычисление безразмерных переменных и значений потенциала в выбранном регионе;
\item расчёт производных потенциала и итоговых свойств;
\item формирование унифицированного результата, включающего свойства и номер региона.
\end{enumerate}

В Region~3 маршрут $pT$ включает шаг определения плотности,
поскольку уравнения заданы в переменных $(\rho,T)$.
Для этого используются обратные зависимости $v(p,T)$,
опубликованные IAPWS в дополнительном релизе \cite{iapws_vpt_region3},
и контроль физической допустимости найденного состояния.

\subsection{Обратные маршруты pH, pS и pX}

Обратные маршруты применяются,
когда состояние задаётся энергетическими параметрами.
К ним относятся $pH$ и $pS$.
В этих постановках нужно восстановить температуру
и далее вычислить полный набор свойств.

Для регионов и подрегионов,
где IF97 предоставляет аналитические обратные формулы,
температура определяется без итераций.
Это снижает вычислительные затраты
в массовых инженерных сценариях.

Если обратная постановка попадает в Region~3,
возникает система нелинейных уравнений относительно $(\rho,T)$.
В проекте для таких 2D-задач применяется решатель
Левенберга--Марквардта с предварительным сеточным поиском
и демпфированием шага.
Этот подход выбран для устойчивой сходимости
в околокритической области.

Маршрут $pX$ используется для двухфазной области.
По заданным давлению $p$ и степени сухости $x$ определяется температура насыщения $T_{\mathrm{sat}}(p)$ и свойства насыщенной жидкости
($x=0$) и насыщенного пара ($x=1$).
Для величин, аддитивных по массе, применяются правила смешения:
\begin{equation}
y = (1-x)\,y_{f} + x\,y_{g},
\end{equation}
где индекс $f$ соответствует насыщенной жидкости, а индекс $g$ --- насыщенному пару.
Удельный объём смеси вычисляется по этой формуле, а плотность определяется как $\rho = 1/v$.
Свойства двухфазной области, включая поведение $c_p$ и $w$,
требуют отдельной интерпретации в прикладных расчётах.

\section{Численные методы расширения модели}

IF97 задаёт основной набор формул и часть обратных зависимостей,
но для полной рабочей области требуются дополнительные численные процедуры.
В проекте они используются как расширение,
сохраняющее физическую основу стандарта.

\subsection{Поиск плотности в Region 3 при задании pT}

В Region~3 давление является функцией плотности и температуры.
При задании состояния в виде $(p, T)$ требуется подобрать $\rho$, удовлетворяющее уравнению
\begin{equation}
F(\rho) = p_{\mathrm{if97}}(\rho, T) - p_{\mathrm{target}} = 0.
\end{equation}

Эта задача решается одномерным поиском корня
с контролем начальных приближений, числа итераций
и принадлежности области Region~3.

Для обратных постановок, приводящих к неизвестным $(\rho, T)$, используется система невязок вида
\begin{equation}
\mathbf{r}(\rho, T) =
\left[
\begin{array}{c}
p_{\mathrm{if97}}(\rho, T) - p_{\mathrm{target}} \\
h_{\mathrm{if97}}(\rho, T) - h_{\mathrm{target}}
\end{array}
\right],
\end{equation}
или аналогично с использованием энтропии.
В околокритической области устойчивость классических ньютоновских шагов
может снижаться из-за плохой обусловленности Якобиана.
Поэтому используется схема Левенберга--Марквардта
с контролем уменьшения нормы $\|\mathbf{r}\|$ и демпфированием шага.
Это повышает надёжность расчёта при резкой нелинейности свойств.

\subsection{Маршрут rhoT методом секущих}

Маршрут $\rho T$ не имеет универсального аналитического выражения в рамках IF97 на всей области применимости.
В проекте он реализован как численное расширение,
основанное на подборе давления $p$, при котором модель IF97 воспроизводит заданную плотность при фиксированной температуре.

Вводится целевая функция
\begin{equation}
F(p) = \rho_{\mathrm{if97}}(p, T) - \rho_{\mathrm{target}}.
\end{equation}
Далее применяется метод секущих:
\begin{equation}
p_{k+1} = p_{k} - F(p_{k})\,\frac{p_{k}-p_{k-1}}{F(p_{k})-F(p_{k-1})}.
\end{equation}

После нахождения давления расчёт продолжается стандартным маршрутом $pT$.
Следовательно, расширение $\rho T$
сохраняет единую физическую основу IF97
и единый формат результата для всех каналов доставки.

\subsection{Особенности rhoT в Region 4}

В двухфазной области (Region~4) температура однозначно задаёт давление насыщения $p_{\mathrm{sat}}(T)$.
Следовательно, при задании $(\rho, T)$ внутри области насыщения давление определяется непосредственно по $T$,
а дополнительной неизвестной является степень сухости $x$.

Для расчёта $x$ удобно использовать удельный объём.
Пусть $v_{\mathrm{target}} = 1/\rho_{\mathrm{target}}$,
а $v_{f}(T)$ и $v_{g}(T)$ --- удельные объёмы насыщенной жидкости и насыщенного пара при температуре $T$.
Тогда степень сухости вычисляется как
\begin{equation}
x = \frac{v_{\mathrm{target}} - v_{f}(T)}{v_{g}(T) - v_{f}(T)}.
\end{equation}
После вычисления $x$ свойства смеси для аддитивных величин определяются правилами смешения.

Вблизи критической точки разность $v_{g}(T)-v_{f}(T)$ уменьшается.
Из-за этого вычисление $x$ становится плохо обусловленным.
В реализации требуется обработка граничных случаев
и контроль диапазона $x \in [0;1]$.

\section{Ограничения применимости и обработка особых случаев}

Корректность расчётов определяется
областью применимости IF97 и устойчивостью численных процедур.
Поэтому в ядре выполняются
валидация входных параметров и явная обработка ошибок.

К типовым ограничениям относятся:
\begin{itemize}
\item допустимые диапазоны давления и температуры в соответствии со стандартом IAPWS-IF97 \cite{iapws_if97_2012};
\item требование $\rho > 0$ для маршрутов, использующих плотность;
\item требование $x \in [0;1]$ для расчётов в двухфазной области;
\item ограничение области метастабильного пара в расширении Region~2 по давлению.
\end{itemize}

К особым случаям относятся вычисления на фазовых границах, в околокритической области и при переходах между регионами.
В этих ситуациях важны:
\begin{itemize}
\item контроль смены региона при итерационном процессе;
\item обработка случаев плохой обусловленности (например, вблизи критической точки);
\item ограничение числа итераций и явная сигнализация о несходимости;
\item единый формат ошибок валидации и вычислений, пригодный для UI и сервисного использования.
\end{itemize}

Перечисленные меры обеспечивают
устойчивую работу модели в инженерных сценариях
и воспроизводимость результатов в \texttt{if97\_core},
FLTK, Yew/WASM+Tauri и gRPC-сервисе.
